\section{Evaluasi}
\label{sec:evaluasi-bab11}

\subsection{Soal Pilihan Ganda}

\begin{enumerate}
    \item Nilai \flag{CF} dan \flag{OF} setelah instruksi
          \lstasm|XOR AL, BL| dijalankan adalah \dots
          \begin{enumerate}
              \item keduanya 0
              \item keduanya 1
              \item \flag{CF} = 1 dan \flag{OF} = 0
              \item bergantung pada nilai \reg{AL} dan \reg{BL}
          \end{enumerate}

    \item Instruksi untuk memadamkan bit ke-3 \reg{AL} tanpa mengubah
          bit lain adalah \dots
          \begin{enumerate}
              \item \lstasm|OR AL, 08h|
              \item \lstasm|AND AL, 08h|
              \item \lstasm|AND AL, 0F7h|
              \item \lstasm|XOR AL, 08h|
          \end{enumerate}

    \item Diketahui \reg{AL} = \bin{1111 1111} yang ditafsirkan
          sebagai $-1$. Instruksi \lstasm|SHR AL, 1| menghasilkan
          nilai \dots
          \begin{enumerate}
              \item $-1$ karena bit tanda dipertahankan
              \item $+127$ karena bit tanda diisi nol
              \item $0$ karena seluruh bit bergeser keluar
              \item $-64$ karena hasilnya dibulatkan ke bawah
          \end{enumerate}

    \item Instruksi yang menyertakan \flag{CF} sebagai bit tambahan
          yang ikut berputar adalah \dots
          \begin{enumerate}
              \item \asm{ROL} dan \asm{ROR}
              \item \asm{RCL} dan \asm{RCR}
              \item \asm{SHL} dan \asm{SHR}
              \item \asm{SAL} dan \asm{SAR}
          \end{enumerate}

    \item Perbedaan hakiki antara \asm{TEST} dan \asm{CMP} terletak
          pada \dots
          \begin{enumerate}
              \item operasi internalnya, yaitu \asm{AND} dan \asm{SUB},
                    meskipun keduanya tidak menyimpan hasil
              \item \asm{TEST} menyimpan hasil sedangkan \asm{CMP}
                    tidak
              \item \asm{TEST} hanya dapat dipakai pada operand 8 bit
              \item keduanya mengubah \flag{CF}, tetapi hanya
                    \asm{CMP} yang mengubah \flag{ZF}
          \end{enumerate}
\end{enumerate}

\subsection{Soal Uraian}

\begin{enumerate}
    \item Jelaskan mengapa \flag{CF} dan \flag{OF} selalu 0 setelah
          operasi logika, sedangkan pada operasi aritmatika keduanya
          dapat bernilai 1. Kaitkan dengan penjalaran \eng{carry}
          antarbit.

    \item Diketahui \reg{AL} = \bin{1100 1010}. Tuliskan hasil
          \lstasm|SHL AL, 1|, \lstasm|SHR AL, 1|, dan
          \lstasm|SAR AL, 1| beserta nilai \flag{CF}, lalu nyatakan
          hasil \asm{SAR} sebagai bilangan bertanda.

    \item Tuliskan program yang menyusun digit puluhan dari variabel
          \code{puluhan} dan digit satuan dari variabel \code{satuan}
          menjadi satu byte BCD, lalu memisahkannya kembali. Jelaskan
          peran \asm{SHL}, \asm{OR}, \asm{AND}, dan \asm{SHR}.

    \item Jelaskan mengapa dekripsi \eng{XOR cipher} dapat memakai
          kunci dan operasi yang sama dengan enkripsi, lalu sebutkan
          dua alasan mengapa cara ini belum memadai sebagai pengamanan
          data.
\end{enumerate}

\subsection{Kisi-kisi Penilaian}

\begin{table}[H]
    \centering
    \caption{Kisi-kisi evaluasi Bab~\ref{chap:logika}.}
    \label{tab:kisi-bab11}
    \small
    \setlength{\tabcolsep}{5pt}
    \begin{tabular}{@{}p{1.5cm}p{5.4cm}p{3.6cm}c@{}}
        \toprule
        \textbf{Sub-CPMK} & \textbf{Indikator yang diukur} &
        \textbf{Bentuk soal} & \textbf{Bobot} \\
        \midrule
        3.2 & Menjelaskan tabel kebenaran dan pengaruh operasi logika
        terhadap \eng{flag} & Pilihan ganda 1; Uraian 1 & 20\% \\
        \addlinespace
        3.2 & Menerapkan manipulasi bit ke-$n$ dan memilih instruksi
        yang tepat & Pilihan ganda 2, 5 & 20\% \\
        \addlinespace
        3.2 & Menganalisis pergeseran logis dan aritmetika pada
        bilangan bertanda & Pilihan ganda 3; Uraian 2 & 25\% \\
        \addlinespace
        3.2 & Menjelaskan rotasi dan peran \flag{CF} & Pilihan ganda
        4 & 10\% \\
        \addlinespace
        3.2 & Menulis program \eng{packing} BCD dan menganalisis
        \eng{XOR cipher} & Uraian 3, 4 & 25\% \\
        \bottomrule
    \end{tabular}
\end{table}

\begin{checklistbox}
\textbf{Checklist Pencapaian Kompetensi Bab~\ref{chap:logika}}

Seorang mahasiswa dinyatakan kompeten pada Sub-CPMK 3.2 apabila:
\begin{itemize}
    \item[\checked] dapat menyusun tabel kebenaran \asm{AND},
          \asm{OR}, \asm{XOR}, dan \asm{NOT};
    \item[\checked] dapat menjelaskan pengaruh operasi logika terhadap
          \flag{CF}, \flag{OF}, \flag{ZF}, \flag{SF}, dan \flag{PF};
    \item[\checked] dapat menuliskan instruksi penyalakan, pemadaman,
          pembalikan, dan pengujian bit ke-$n$;
    \item[\checked] dapat membedakan \asm{TEST} dan \asm{CMP} beserta
          pemakaiannya;
    \item[\checked] dapat menentukan hasil \asm{SHL}, \asm{SHR}, dan
          \asm{SAR} pada bilangan bertanda disertai \flag{CF};
    \item[\checked] dapat menjelaskan perbedaan \asm{ROL}/\asm{ROR}
          dengan \asm{RCL}/\asm{RCR};
    \item[\checked] berhasil menulis dan menjalankan program paritas
          serta \eng{packing} BCD.
\end{itemize}
\end{checklistbox}
